% ==============================================================================
% Fichier : front/02_mode_emploi.tex
% Titre   : Mode d'emploi — GCI
% Auteur  : MINKA MI NGUIDJOÏ Thierry Emmanuel
% Version : 1.0 -- octobre 2026
% ==============================================================================

\chapter*{Mode d'emploi de ce cours}
\addcontentsline{toc}{chapter}{Mode d'emploi}
\markboth{Mode d'emploi}{Mode d'emploi}

\section*{Structure générale}

Le cours est organisé en \textbf{6 parties précédées d'un prologue}, soit
21~chapitres numérotés de façon continue. Il accompagne une unité
d'enseignement de 30~heures en présentiel, réparties en quinze séances de
deux heures qui couvrent l'ensemble des parties, du prologue à la
partie~VI. Les ateliers de questions-réponses par chapitre et les
laboratoires de NDAP SUKULU sont des travaux à faire à la maison, hors de
ces 30~heures. Les parties~V et~VI servent en outre de point de départ aux
mémoires.

\begin{tableaucours}{lp{3.8cm}p{3.8cm}c}{%
    \tableauHeader{Partie} &
    \tableauHeader{Thème} &
    \tableauHeader{Compétences visées} &
    \tableauHeader{Niveau}
}
0 & Prologue opérationnel & Distinguer une cause technique d'une cause de gouvernance & \niveaubase \\
I & Fondements de la gouvernance des SI & Définir, situer et justifier la gouvernance des SI & \niveaubase \\
II & COBIT 2019 & Concevoir et évaluer un système de gouvernance & \niveaubase\,--\,\niveaumoyen \\
III & Cadre normatif et juridique & Identifier les obligations et auditer la conformité & \niveaubase \\
IV & ITIL 4 et processus & Organiser et mesurer la gestion des services & \niveaumoyen\,--\,\niveauexpert \\
V & Défis avancés & Gouverner l'agilité, le nuage, les données et l'IA & \niveauexpert \\
VI & Recherche & Lire et produire la recherche sur la gouvernance des SI & \niveauexpert \\
\end{tableaucours}

\section*{Les environnements visuels}

\begin{boxdef}
\textbf{Définition}. Introduit un terme technique avec sa définition formelle.
Tous les termes définis sont indexés en fin d'ouvrage.
\end{boxdef}

\begin{boxobjectifs}
Objectifs du chapitre. Liste les compétences mesurables attendues à l'issue
du chapitre.
\end{boxobjectifs}

\begin{boxscenario}{Titre du scénario}
Scénario opérationnel. Présente un cas fictif ancré dans la réalité
camerounaise ou africaine, qui revient d'un chapitre à l'autre (opérations
SANAGA, KRIBI, MEFOU et NOUN).
\end{boxscenario}

\begin{boxmetier}{DSI}
Application métier. Traduit le concept académique en démarche
opérationnelle pour l'un des cinq métiers visés : DSI, Responsable GRC,
Auditeur SI, Responsable des services IT et RSSI.
\end{boxmetier}

\begin{boxprocedure}
\libelleProcedure. Indique ce qui doit être documenté et prouvé pour qu'une
pratique soit auditable : la pièce attendue, son responsable, sa fréquence.
\end{boxprocedure}

\begin{boxobjet}{Code et nom}
\libelleObjet. Présente un objectif COBIT ou une pratique ITIL sous une forme
constante : finalité, pratiques, entrées et sorties, matrice RACI,
indicateurs.
\end{boxobjet}

\begin{boxjuris}
Jurisprudence \& Droit applicable. Référence les textes de loi et articles
pertinents, avec accent sur le droit camerounais et africain.
\end{boxjuris}

\begin{boxalert}
Attention. Signale une erreur fréquente de gouvernance, de conformité ou de
mise en œuvre.
\end{boxalert}

\begin{boxresume}
À retenir. Synthèse des points essentiels du chapitre.
\end{boxresume}

\begin{boxexo}
Exercice. Travaux gradués : \niveaubase\ fondamental,
\niveaumoyen\ intermédiaire, \niveauexpert\ expert.
\end{boxexo}

\section*{Conventions typographiques}

\begin{itemize}
    \item Les \termecle{termes clés} apparaissent en \textbf{gras bleu marine}
          et sont indexés automatiquement.
    \item Les \sigle{ACRONYMES} utilisent des petites capitales.
    \item Les références aux normes utilisent la commande \norme{ISO~XXXX}.
    \item Les références aux lois camerounaises utilisent \loicam{AAAA/NNN}.
    \item Les niveaux de capacité COBIT s'écrivent \capacite{2}{DSS02}.
\end{itemize}

\section*{Matériel d'accompagnement}

\begin{itemize}
    \item \textbf{NDAP SUKULU}, code de cours \codeUE~: une fiche
          d'auto-évaluation par chapitre et dix laboratoires pratiques de
          60~minutes.
    \item \textbf{Dossiers de cas}~: pièces des quatre organisations fictives
          (tableurs, extraits de tickets, contrats, procès-verbaux).
    \item \textbf{Modèles de documents} (Annexe~\ref{ann:templates})~: charte de gouvernance,
          matrice RACI, accord de niveau de service, grille d'évaluation de
          capacité.
\end{itemize}
